\section{Series, hypergeometric form, branch points, and the differential resolvent}\label{sec:series}

\subsection{The Taylor coefficients by finite Lagrange inversion}

Put
\[
 y(a)=-\BR(a).
\]
Then \(y(0)=0\) and
\begin{equation}\label{eq:y-fixed}
 y+y^5=a,
 \qquad
 y=a(1+y^4)^{-1}.
\end{equation}
In the truncated ring \(\Q[a]/(a^{N+1})\), the inverse
\((1+y^4)^{-1}\) is a finite geometric sum.  Theorem~\ref{thm:finite-lagrange}
therefore applies with \(\Phi(w)=(1+w^4)^{-1}\).

\begin{theorem}[Bring coefficient formula]\label{thm:bring-coeff}
For every \(N\), the unique truncated solution of~\eqref{eq:y-fixed} is
\begin{equation}\label{eq:bring-series-y}
 y(a)=\sum_{4k+1\le N}
 (-1)^k\frac{1}{4k+1}\binom{5k}{k}a^{4k+1}
 \pmod {a^{N+1}}.
\end{equation}
Consequently
\begin{equation}\label{eq:bring-series}
 \BR(a)=\sum_{k\ge0}
 (-1)^{k+1}\frac{1}{4k+1}\binom{5k}{k}a^{4k+1}
 =-a+a^5-5a^9+35a^{13}-285a^{17}+\cdots .
\end{equation}
The assertion through every requested degree is a theorem of \ipra.
\end{theorem}

\begin{proof}
Finite Lagrange inversion gives
\[
 [a^n]y=\frac1n[w^{n-1}](1+w^4)^{-n}.
\]
The finite negative-binomial identity is
\[
 (1+w^4)^{-n}
 =\sum_{j=0}^{\lfloor N/4\rfloor}
   (-1)^j\binom{n+j-1}{j}w^{4j}
   \pmod {w^{N+1}}.
\]
Thus the coefficient vanishes unless \(n-1=4k\).  For
\(n=4k+1\),
\[
 \binom{n+k-1}{k}=\binom{5k}{k},
\]
which gives~\eqref{eq:bring-series-y}.  In the object theory one keeps the
integer equation
\((4k+1)[a^{4k+1}]y=(-1)^k\binom{5k}{k}\), so no unguarded division is used.
\end{proof}

The archive's exact checker forms the right side through degree 65 and verifies
that \(y^5+y-a\) has no coefficient through that degree.  This finite check is
not the proof of the uniform theorem, but it is a useful independent regression
test.

\subsection{The hypergeometric identity is a coefficient recurrence}

Let
\[
 c_k=\frac{1}{4k+1}\binom{5k}{k}
     =\frac{(5k)!}{k!(4k+1)!}.
\]
A direct cancellation of factorials gives
\begin{equation}\label{eq:bring-ratio}
 \frac{c_{k+1}}{c_k}
 =\frac{3125}{256}
 \frac{(k+\frac15)(k+\frac25)(k+\frac35)(k+\frac45)}
 {(k+\frac12)(k+\frac34)(k+\frac54)(k+1)}.
\end{equation}

\begin{theorem}[Hypergeometric form]\label{thm:bring-hypergeom}
On its disk of convergence,
\begin{equation}\label{eq:bring-hypergeom}
 \BR(a)=-a\,{}_4F_3\!\left(
 \begin{matrix}\frac15,\frac25,\frac35,\frac45\\
               \frac12,\frac34,\frac54
 \end{matrix};-\frac{3125}{256}a^4\right).
\end{equation}
For every truncation order, this is an exact identity of rational coefficient
vectors and is therefore an \ipra{} theorem.
\end{theorem}

\begin{proof}
The coefficient of \(a^{4k+1}\) on the right is
\[
 -\frac{(\frac15)_k(\frac25)_k(\frac35)_k(\frac45)_k}
 {(\frac12)_k(\frac34)_k(\frac54)_k k!}
 \left(-\frac{3125}{256}\right)^k.
\]
At \(k=0\) it equals \(-1\).  Its ratio to the preceding coefficient is the
negative of~\eqref{eq:bring-ratio}, exactly matching
\((-1)^{k+1}c_k\).  Equality follows by induction on \(k\).  After clearing the
fixed denominators \(2,3,4,5\), the recurrence is an integer polynomial
identity, so quantifier-free induction proves the finite vector equality.
\end{proof}

This proof is preferable in the present setting to invoking a general theory
of hypergeometric functions: the named special function is simply the regular
sequence of its rational partial sums plus the tail modulus established next.

\subsection{Exact convergence radius and boundary convergence}

Taking \(k\to\infty\) in~\eqref{eq:bring-ratio} gives
\[
 \lim_{k\to\infty}\frac{c_{k+1}}{c_k}=\frac{3125}{256}.
\]
The limit is constructive.  Subtracting \(3125/256\) from the rational function
in~\eqref{eq:bring-ratio} gives a quotient whose numerator has degree at most
three and whose positive denominator has degree four.  A fixed integer
\(C\) therefore satisfies
\begin{equation}\label{eq:ratio-modulus}
 \left|\frac{c_{k+1}}{c_k}-\frac{3125}{256}\right|
 \le\frac{C}{k+1}.
\end{equation}
The inequality is checked after multiplication by the positive denominator,
and is thus available in \ipra{} with an explicit cutoff.

\begin{theorem}[Radius of the Bring series]\label{thm:bring-radius}
The radius of convergence of~\eqref{eq:bring-series} is
\begin{equation}\label{eq:bring-radius}
 R=\left(\frac{256}{3125}\right)^{1/4}
   =\frac{4}{5\sqrt[4]{5}}.
\end{equation}
On every named closed disk \(|a|\le r<R\), the series has a
primitive-recursive geometric tail modulus.  If \(|a|>R\), its terms do not
tend to zero.  The series also converges absolutely on \(|a|=R\), with a
primitive-recursive tail modulus of order \(N^{-1/2}\).
\end{theorem}

\begin{proof}
Regard the series after its leading factor \(a\) as a series in \(t=a^4\).
Given rational \(r<R\), choose a rational \(\theta<1\) strictly above
\((3125/256)r^4\).  Estimate~\eqref{eq:ratio-modulus} supplies an explicit
\(K\) after which consecutive absolute terms have ratio at most \(\theta\),
so a geometric tail bound follows.  For \(|a|>R\), choose a rational ratio
strictly greater than one below the eventual term ratio; the terms cannot tend
to zero.

At the boundary, the finite Wallis--Stirling inequalities used in the ProveIt
Lagrange-inversion packet give a rational constant \(C_0\) such that
\[
 c_k\left(\frac{256}{3125}\right)^k
 \le \frac{C_0}{(k+1)^{3/2}}.
\]
The proof uses only factorial products, rational upper and lower bounds for
\(\log(1+u)\), and dyadic grouping.  The tail of
\(\sum(k+1)^{-3/2}\) is bounded by a constant times \(N^{-1/2}\), again by
dyadic grouping.  Every cutoff and rational constant is primitive recursive.
\end{proof}

\subsection{The four branch points and the discriminant}

A multiple root of \(x^5+x+a\) satisfies
\[
 x^5+x+a=0,
 \qquad 5x^4+1=0.
\]
The second equation gives \(x^4=-1/5\), and hence
\(x^5=-x/5\) and \(a=-4x/5\).  Therefore
\begin{equation}\label{eq:branch-equation}
 3125a^4+256=0.
\end{equation}
Conversely every solution of~\eqref{eq:branch-equation}, with
\(x=-5a/4\), satisfies both equations.

\begin{theorem}[Trinomial discriminant]\label{thm:trinomial-discriminant}
For indeterminates \(p,q\),
\begin{equation}\label{eq:trinomial-discriminant}
 \Disc_X(X^5+pX+q)=256p^5+3125q^4.
\end{equation}
In particular, the four critical values of the Bring projection are the four
roots of~\eqref{eq:branch-equation}; all have modulus \(R\) from
\eqref{eq:bring-radius}.
\end{theorem}

\begin{proof}
The discriminant is the resultant of
\(X^5+pX+q\) and \(5X^4+p\), since
\((-1)^{5\cdot4/2}=1\).  The determinant of the fixed \(9\times9\) Sylvester
matrix normalizes to the right side of~\eqref{eq:trinomial-discriminant}.
Equivalently, eliminate \(X\): a common root has
\(X^4=-p/5\) and
\(q=-4pX/5\); raising the latter equation to the fourth power and substituting
the former yields the displayed relation.  The converse is checked after
adjoining a fourth root.  Both versions are finite polynomial arguments.
\end{proof}

At a critical root \(x_0\), one has \(x_0\ne0\) and
\(P_{xx}(x_0)=20x_0^3\ne0\).  Thus the local inverse has two Puiseux branches
with leading term proportional to \((a-a_0)^{1/2}\).  Analytic continuation
around a small rational polygonal loop enclosing \(a_0\) interchanges these
two branches.  A single-valued analytic branch on the whole complex plane is
therefore impossible, and any domain for a branch agreeing with the real Bring
radical must introduce cuts.  The local square-root construction is finite in
every truncation and fits the ramified Lagrange theorem of the reused packet;
the loop interchange can be certified by a finite winding-number computation.

The critical values lie exactly on the circle \(|a|=R\).  Hence the branch
point calculation independently proves that the Taylor disk at zero cannot be
larger than the radius found by the coefficient ratio.

\subsection{The Cockle--Harley differential resolvent}

The differential-resolvent convention uses a branch \(t=t(\rho)\) satisfying
\begin{equation}\label{eq:minus-bring}
 t^5-t+\rho=0.
\end{equation}
Write \(D=5t^4-1\).  Implicit differentiation gives
\(t'=-D^{-1}\).  Repeatedly applying the total derivative
\[
 \mathcal D=\frac{\partial}{\partial\rho}
            -\frac1D\frac{\partial}{\partial t}
\]
expresses \(t'',t^{(3)},t^{(4)}\) as rational functions of \(t\), with the
fourth derivative having denominator \(D^7\).

\begin{theorem}[Fourth-order differential resolvent]\label{thm:differential-resolvent}
Every simple branch of~\eqref{eq:minus-bring} satisfies
\begin{equation}\label{eq:cockle-harley}
 \frac{256-3125\rho^4}{1155}\,t^{(4)}
 -\frac{6250\rho^3}{231}\,t^{(3)}
 -\frac{4875\rho^2}{77}\,t''
 -\frac{2125\rho}{77}\,t'
 +t=0.
\end{equation}
For every fixed Taylor order this is an \ipra{} theorem about rational
coefficient vectors; on a named disk avoiding the critical values it is an
effective analytic theorem.
\end{theorem}

\begin{proof}
Substitute the four rational derivative expressions into the left side,
multiply by \(1155D^7\), and replace \(\rho\) by \(t-t^5\) using
\eqref{eq:minus-bring}.  Polynomial normalization gives zero.  This is the
exact certificate implemented by \code{check\_differential\_resolvent} in the
archive.  In \ipra{}, the same calculation is an integer ring identity under
the guarded hypothesis \(D\ne0\).  At a simple branch the guard follows from
the nonvanishing derivative.
\end{proof}

At \(\rho=0\), the five roots are \(0,1,-1,i,-i\), all simple.  The branch at
zero is the hypergeometric branch of Theorem~\ref{thm:bring-hypergeom} after the
appropriate sign change.  For a fourth root of unity \(\zeta\), write
\(t=\zeta+u\); then
\[
 u=-\rho\bigl(4+10\zeta^3u+10\zeta^2u^2+5\zeta u^3+u^4\bigr)^{-1}.
\]
Finite fixed-point iteration constructs the other four Taylor branches.  They
all satisfy~\eqref{eq:cockle-harley}; their sum is zero by the missing quartic
coefficient, explaining why a fourth-order equation suffices for five
algebraic branches.  This gives a constructive derivation of all five local
roots even when one does not adopt the particular hypergeometric basis printed
by MathWorld.
